% !TeX root = ../queens.tex
\section{Introduction}
Let $\N=\{0,1,2,\ldots\}$. Set $q_0=0$ and, for
$n\ge1$, let $q_n$ be the least $y\in\N$ such that
\begin{equation}\label{eq:greedy}
 y\ne q_i,\qquad y-n\ne q_i-i,\qquad y+n\ne q_i+i
 \qquad(0\le i<n);
\end{equation}
that is, $(n,y)$ does not share a row, diagonal, or antidiagonal
with any previous queen.
Each step excludes only finitely many rows, so $q_n$ is defined
for every column $n$, and these queens are pairwise nonattacking. The sequence begins
\[
 (q_0,q_1,q_2,\dots)=(0,2,4,1,3,8,10,12,14,5,7,18,6,21,9,\dots);
\]
it is the \emph{greedy queens permutation of $\N$} (we prove that it is a permutation in Lemma~\ref{lem:surjective}).
Call a queen \textcolor{queensUpper}{\emph{upper}} if $q_n>n$ and
\textcolor{queensLower}{\emph{lower}} if $q_n<n$.
For a square $(x,y)$, its diagonal is indexed by $y-x$ and its antidiagonal by $y+x$.

\begin{figure}[!ht]
\centering
\definecolor{queensSquare}{HTML}{F2F3F5}
\definecolor{queensGuide}{HTML}{9CA6B8}
\definecolor{queensFocus}{HTML}{E7EBFF}
\definecolor{queensAnti}{HTML}{5865A8}
% Scale the three panels together to keep the figure on the first page.
\scalebox{0.90}{%
% A compact group with closely fitted columns centers the visible artwork.
\begin{minipage}{0.84\linewidth}
\begin{minipage}[t]{0.50\linewidth}
\vspace{0pt}
\centering
\setboardfontsize{9.5pt}
\begin{tikzpicture}[x=0.3cm,y=0.3cm,
  coordinate label/.style={font=\sffamily\scriptsize,text=black!60,inner sep=1pt},
  queen/.style={inner sep=0pt,text=black}]
  % The greedy decision in column 20; (0,0) is at lower left.
  \foreach \x in {0,...,20} {
    \foreach \y in {0,...,30} {
      \pgfmathtruncatemacro{\parity}{mod(\x+\y,2)}
      \ifnum\parity=0
        \fill[queensSquare] (\x-0.5,\y-0.5) rectangle (\x+0.5,\y+0.5);
      \fi
    }
  }
  \fill[queensFocus] (19.5,-0.5) rectangle (20.5,30.5);
  \draw[queensGuide,dash pattern=on 2.2pt off 2.2pt,line width=0.45pt]
    (-0.5,-0.5) -- (20.5,20.5);
  % The first 20 queens: teal above y=n, burnt orange below; the origin is black.
  \node[queen] at (0,0) {\BlackQueenOnWhite};
  \foreach \x/\y in {1/2,2/4,5/8,6/10,7/12,8/14,11/18,13/21,
                      15/24,16/26,17/28,18/30} {
    \node[queen,text=queensUpper] at (\x,\y) {\BlackQueenOnWhite};
  }
  \foreach \x/\y in {3/1,4/3,9/5,10/7,12/6,14/9,19/11} {
    \node[queen,text=queensLower] at (\x,\y) {\BlackQueenOnWhite};
  }
  % Row 13 is the only unattacked displayed square in column 20.
  \foreach \y in {0,...,12,14,15,...,30} {
    \draw[queensGuide,line width=0.35pt]
      (19.88,\y-0.12) -- (20.12,\y+0.12)
      (19.88,\y+0.12) -- (20.12,\y-0.12);
  }
  \node[queen,text=queensLower] at (20,13) {\BlackQueenOnWhite};
  \foreach \i in {0,2,...,20} {
    \node[coordinate label,anchor=north] at (\i,-0.65) {\i};
  }
  \foreach \i in {0,5,...,30} {
    \node[coordinate label,anchor=east] at (-0.65,\i) {\i};
  }
\end{tikzpicture}
\end{minipage}\hfill
\begin{minipage}[t]{0.40\linewidth}
\vspace{0pt}
\centering
% Each small square contains its index, so equal labels trace one diagonal.
\begin{tikzpicture}[x=0.6cm,y=0.6cm,
  index/.style={font=\small,text=black!40,inner sep=0pt},
  coordinate label/.style={font=\sffamily\scriptsize,text=black!60,inner sep=1pt}]
  \node[font=\small\bfseries] at (2,5.15) {Diagonals: $y-x$};
  \foreach \x in {0,...,4} {
    \foreach \y in {0,...,4} {
      \pgfmathtruncatemacro{\parity}{mod(\x+\y,2)}
      \ifnum\parity=0
        \fill[queensSquare] (\x-0.5,\y-0.5) rectangle (\x+0.5,\y+0.5);
      \fi
      \pgfmathtruncatemacro{\diagonalindex}{\y-\x}
      \ifnum\diagonalindex=2
        \fill[queensUpper!15] (\x-0.5,\y-0.5) rectangle (\x+0.5,\y+0.5);
        \node[index,text=queensUpper,font=\small\bfseries] at (\x,\y) {$2$};
      \else\ifnum\diagonalindex=-2
        \fill[queensLower!17] (\x-0.5,\y-0.5) rectangle (\x+0.5,\y+0.5);
        \node[index,text=queensLower,font=\small\bfseries] at (\x,\y) {$-2$};
      \else
        \node[index] at (\x,\y) {$\diagonalindex$};
      \fi\fi
    }
  }
  \foreach \i in {0,...,4} {
    \node[coordinate label,anchor=north] at (\i,-0.65) {\i};
    \node[coordinate label,anchor=east] at (-0.65,\i) {\i};
  }
  \draw[black,line width=0.5pt]
    (-0.5,4.5) -- (-0.5,-0.5) -- (4.5,-0.5);
  \node[coordinate label,text=black,anchor=west,inner sep=3pt]
    at (4.5,-0.5) {$x$};
  \node[coordinate label,text=black] at (-0.9,4.6) {$y$};
  \node[font=\footnotesize] at (2,-1.55)
    {\textcolor{queensUpper}{2nd upper diagonal}:\quad $y-x=2$};
  \node[font=\footnotesize] at (2,-2.15)
    {\textcolor{queensLower}{2nd lower diagonal}:\quad $y-x=-2$};
\end{tikzpicture}

\vspace{0.35cm}

\begin{tikzpicture}[x=0.6cm,y=0.6cm,
  index/.style={font=\small,text=black!40,inner sep=0pt},
  coordinate label/.style={font=\sffamily\scriptsize,text=black!60,inner sep=1pt}]
  \node[font=\small\bfseries] at (2,5.15) {Antidiagonals: $y+x$};
  \foreach \x in {0,...,4} {
    \foreach \y in {0,...,4} {
      \pgfmathtruncatemacro{\parity}{mod(\x+\y,2)}
      \ifnum\parity=0
        \fill[queensSquare] (\x-0.5,\y-0.5) rectangle (\x+0.5,\y+0.5);
      \fi
      \pgfmathtruncatemacro{\antidiagonalindex}{\y+\x}
      \ifnum\antidiagonalindex=3
        \fill[queensAnti!16] (\x-0.5,\y-0.5) rectangle (\x+0.5,\y+0.5);
        \node[index,text=queensAnti,font=\small\bfseries] at (\x,\y) {$3$};
      \else
        \node[index] at (\x,\y) {$\antidiagonalindex$};
      \fi
    }
  }
  \foreach \i in {0,...,4} {
    \node[coordinate label,anchor=north] at (\i,-0.65) {\i};
    \node[coordinate label,anchor=east] at (-0.65,\i) {\i};
  }
  \draw[black,line width=0.5pt]
    (-0.5,4.5) -- (-0.5,-0.5) -- (4.5,-0.5);
  \node[coordinate label,text=black,anchor=west,inner sep=3pt]
    at (4.5,-0.5) {$x$};
  \node[coordinate label,text=black] at (-0.9,4.6) {$y$};
  \node[font=\footnotesize] at (2,-1.55)
    {\textcolor{queensAnti}{3rd antidiagonal}:\quad $y+x=3$};
\end{tikzpicture}
\end{minipage}
\end{minipage}%
}
\caption{The greedy step in column $20$ (left), diagonal labels $y-x$
(upper right), and antidiagonal labels $y+x$ (lower right).}
\label{fig:infinite-board}
\end{figure}

Antti Karttunen introduced the positive-integer version of this sequence
in 2001 as OEIS \href{https://oeis.org/A065188}{A065188}; our $0$-indexing gives
\href{https://oeis.org/A275895}{A275895} \cite{OEIS}.
Dekking, Shallit, and Sloane studied a construction that scans the
board along successive antidiagonals rather than columns; they showed that it produces the same set of
queen positions as our column-by-column rule \cite[pp.~24--25]{DSS}.
They conjectured that upper and lower queens stay within bounded
vertical distances of the lines $y=\phi x$ and $y=\phi^{-1}x$ respectively
\cite[Conjecture~25]{DSS}, where
\[
 \phi=\frac{1+\sqrt5}{2}\approx1.618034,\qquad \phi^{-1}=\phi-1\approx0.618034.
\]

\begin{theorem}\label{thm:main}
For every $n\ge1$,
\[
 \left\{
 \begin{aligned}
 \left|q_n-n\phi\right|&<\frac{5}{\phi}\approx3.09017&&\text{if }q_n>n,\\[4pt]
 \Bigl|q_n-\frac{n}{\phi}\Bigr|&<4+\frac{5}{\phi}\approx7.09017&&\text{if }q_n<n.
 \end{aligned}
 \right.
\]
\end{theorem}

\begin{figure}[!ht]
\centering
\begin{minipage}[c]{0.36\linewidth}
\centering
\begin{tikzpicture}[x=0.031cm,y=0.031cm,
  tick label/.style={font=\small,text=black!65,inner sep=2pt},
  line label/.style={font=\small,anchor=west,inner sep=4pt}]
  % Queen positions from the greedy rule, for columns 0 through 100.
  \draw[black,line width=0.5pt] (0,170) -- (0,0) -- (105,0);
  \foreach \x in {0,20,40,60,80,100} {
    \draw[black,line width=0.4pt] (\x,0) -- (\x,-2);
    \node[tick label,anchor=north] at (\x,-3) {\x};
  }
  \foreach \y in {0,50,100,150} {
    \draw[black,line width=0.4pt] (0,\y) -- (-2,\y);
    \node[tick label,anchor=east] at (-3,\y) {\y};
  }
  \node[font=\small,anchor=west,inner sep=3pt] at (105,0) {$x$};
  \node[font=\small,anchor=south,inner sep=3pt] at (0,170) {$y$};
  % Draw the reference lines first, so the actual queen positions stay visible.
  \draw[queensUpper!45,line width=0.5pt]
    (0,0) -- (100,{100*(1+sqrt(5))/2});
  \draw[queensLower!45,line width=0.5pt]
    (0,0) -- (100,{200/(1+sqrt(5))});
  \node[line label,text=queensUpper] at (102,{100*(1+sqrt(5))/2})
    {$y=x\phi$};
  \node[line label,text=queensLower] at (102,{200/(1+sqrt(5))})
    {$y=\dfrac{x}{\phi}$};
  \foreach \n/\qn in {
    1/2,2/4,3/1,4/3,5/8,6/10,7/12,8/14,9/5,10/7,
    11/18,12/6,13/21,14/9,15/24,16/26,17/28,18/30,19/11,20/13,
    21/34,22/36,23/38,24/40,25/15,26/17,27/44,28/16,29/47,30/19,
    31/50,32/52,33/20,34/55,35/57,36/59,37/22,38/62,39/23,40/65,
    41/27,42/25,43/69,44/71,45/73,46/75,47/77,48/29,49/31,50/81,
    51/83,52/85,53/32,54/88,55/33,56/91,57/37,58/35,59/95,60/97,
    61/99,62/101,63/39,64/104,65/106,66/41,67/109,68/42,69/112,70/43,
    71/115,72/117,73/119,74/45,75/122,76/46,77/49,78/126,79/48,80/129,
    81/131,82/133,83/135,84/51,85/53,86/139,87/141,88/143,89/54,90/56,
    91/147,92/149,93/151,94/58,95/154,96/156,97/158,98/60,99/161,100/61} {
    \ifnum\qn>\n
      \filldraw[fill=queensUpper,draw=white,line width=0.3pt]
        (\n,\qn) circle[radius=1.1pt];
    \else
      \filldraw[fill=queensLower,draw=white,line width=0.3pt]
        (\n,\qn) circle[radius=1.1pt];
    \fi
  }
  \fill[black] (0,0) circle[radius=1.1pt];
\end{tikzpicture}
\end{minipage}\hfill
\begin{minipage}[c]{0.60\linewidth}
Larsson and W\"astlund \cite{LW} obtained analogous golden-ratio
bounds for Maharaja Nim, where the queen can also make knight moves. Dekking, Shallit, and Sloane
\cite[p.~26]{DSS} asked whether their method could be adapted to the
single-quadrant greedy-queens problem. Our proof shares the strategy of \cite{LW}: we use a finite symbolic
model to show that the $j$th lower queen lies within a fixed distance
of the $j$th lower diagonal, then deduce the golden-ratio bounds
from counting identities. 
We use different methods to establish boundedness however, since the single-quadrant greedy-queens problem is asymmetric, unlike Maharaja Nim.

\end{minipage}
\caption{Queen positions $(n,q_n)$ for $0\le n\le100$ on equally scaled
axes, with the reference lines $y=x\phi$ and $y=x/\phi$.}
\label{fig:golden-lines}
\end{figure}

Carter \cite[Section~5]{Carter} gives a related local description of the
single-quadrant problem and proposes a two-string rewriting approach,
whose completeness is left unproved. Our construction uses the same
separation of upper and lower attacks, but encodes upper-queen information
by a \emph{history graph}. We establish the required bounds by
finite-state verification and an induction connecting the verified
transitions to the actual greedy process.
Like Larsson and W\"astlund \cite{LW}, our proof reduces Theorem~\ref{thm:main} to a bound on lower diagonals: namely, it suffices to prove that $|d_j-j|\le4$ (Lemma~\ref{lem:estimates}) for all $j\ge1$, where $d_j$ is such that
the $j$th lower queen lies on the $d_j$th lower diagonal. 

\begin{remark}[Game interpretation]
Consider a single moving queen on an otherwise empty $\N^2$ board.
Two players alternate moves of the form
\[
 (x,y)\longmapsto(x-t,y),\quad(x,y-t),\quad
 (x-t,y-t),\quad(x-t,y+t),
\]
where $t\ge1$ is an integer and the destination has nonnegative
coordinates. The player unable to move loses. Every move decreases
$2x+y$, so play terminates. The losing positions are exactly $(x,q_x)$:
no move joins two chosen squares, while every other square can reach
one. Indeed, if $y>q_x$, move down to $(x,q_x)$; if $y<q_x$, the greedy
rule supplies an attacking queen in an earlier column.
Thus the chosen squares are the $\mathcal P$-positions, or positions
of Sprague--Grundy value zero \cite[Sections~7--8]{DSS}.
Writing $G$ for the Sprague--Grundy function, $G(x,y)=0$ exactly when $y=q_x$.
Theorem~\ref{thm:main} bounds the location of this zero set, without
determining the positive values.
Deleting the antidiagonal move gives Wythoff Nim, whose losing positions
have an exact golden-ratio description \cite[Section~1]{LW}.
\end{remark}

\medskip
\noindent\textbf{Paper organization.} 
Section~\ref{sec:rows} proves some useful lemmas about greedy queen positions. 
Section~\ref{sec:contraction} uses these lemmas to turn the bound on lower 
diagonals into a recurrence for the number of upper queens up to the $n$th column. 
A contraction argument then yields the stated golden-ratio estimates.
Section~\ref{sec:local} develops an encoding for local states of the board.
Section~\ref{sec:sufficiency} proves that these local states suffice for determining the position of greedy queens on the board.
Section~\ref{sec:finite} verifies closure and the diagonal discrepancy bound for
a finite collection of states by computer, proves by induction
that the actual greedy process always stays within this collection, and
sharpens Theorem~\ref{thm:main}'s constants.
Section~\ref{sec:generation} uses our results to give a sequential
algorithm for generating $q_0,\ldots,q_n$ in $O(n)$ time using
$O(\log n)$ words of working memory.
Appendix~\ref{app:artifacts} discusses how to reproduce our results with a computer.

\medskip
\noindent\textbf{Notation.} For real $x,y$, write $[x\dts y]=\{k\in\mathbb Z: x\le k\le y\}$. For a proposition $P$, the \emph{Iverson bracket}
$[P]$ equals $1$ if $P$ is true and $0$ otherwise.